% Potenzen, Zehnerpotenzen und Quadratwurzeln · Quadratwurzeln · Prüfungs-Fokus MSA – bank/potenzen-wurzeln/e3.jsonl (zusammenbau v0.6, Prüfungs-Fokus)
\einheitenkopf[e3]{Einheit 3 · Quadratwurzeln}
\zweigzeile{ab Kl. 7, je nach Buch bis Kl. 9 (am Gymnasium ab Kl. 8) · P10 ×12}
\setcounter{aufgabe}{0}

% Anlauf (ohne Prüfkennung)
\begin{aufgabe}{Quadratwurzeln – Anlauf}
\begin{teile}
\teil $\sqrt{121}$ – geht die Wurzel im Kopf auf?\\ \kreuz{ja, Quadratzahl} \\ \kreuz{nein, Näherungswert}
\teil $\sqrt{49}$ – wie viel? Mache die Probe. \leerfeld , Probe: \leerfeld
\teil $\sqrt{(-8)^2}$ – wie viel? \leerfeld
\end{teile}
\end{aufgabe}

\begin{aufgabe}{Quadratwurzeln – Prüfungsaufgaben}
\begin{teile}
\teil Ein quadratischer Teppich hat einen Flächeninhalt von $6{,}25\,\text{m}^2$. Wie lang ist eine Seite? \hfill (P20) \\ \leerfeld[m]
\teil Ein quadratischer Platz hat einen Flächeninhalt von $1\,600\,\text{m}^2$. Wie lang ist eine Seite? \hfill (P20) \\ \leerfeld[m]
\teil Eine quadratische Fliese hat einen Flächeninhalt von $441\,\text{cm}^2$. Wie lang ist eine Seite? \hfill (P20) \\ \leerfeld[cm]
\end{teile}
\end{aufgabe}

\begin{aufgabe}{Quadratwurzeln – Prüfungsaufgaben – weiter}
\begin{teile}
\steil Welche Aussage ist wahr? Kreuze an. \hfill (P15*)\\ \kreuz{$\sqrt{(-7)^2} = -7$} \\ \kreuz{$\sqrt{(-7)^2} = 7$} \\ \kreuz{$\sqrt{(-7)^2}$ ist nicht definiert}
\steil Welche Aussage ist wahr? Kreuze an. \hfill (P15*)\\ \kreuz{$\sqrt{(-13)^2} = -13$} \\ \kreuz{$\sqrt{(-13)^2} = 13$} \\ \kreuz{$\sqrt{(-13)^2}$ ist nicht definiert}
\teil Welche Aussage ist wahr? Kreuze an. \hfill (P15)\\ \kreuz{$2{,}25 < \frac{9}{4}$} \\ \kreuz{$\frac{7}{3} > \frac{9}{4}$} \\ \kreuz{$\sqrt{5} > \frac{9}{4}$}
\teil Welche Aussage ist wahr? Kreuze an. \hfill (P15)\\ \kreuz{$1{,}75 > \frac{7}{4}$} \\ \kreuz{$\frac{9}{5} < \frac{7}{4}$} \\ \kreuz{$\sqrt{3} < \frac{7}{4}$}
\teil $\sqrt{3}$; $-\frac{1}{4}$; $1{,}7$; $-0{,}26$ – aufsteigend geordnet? \hfill (P14)
\teil $\sqrt{11}$; $-\frac{3}{5}$; $3{,}3$; $-0{,}62$ – aufsteigend geordnet? \hfill (P14)
\teil Eine Leiter ist $7{,}5$ m lang. Ihr Fuß steht $2{,}1$ m von der Hauswand entfernt. Wie hoch reicht sie an der Wand? \hfill (P24) \\ \leerfeld[m]
\teil Eine Rampe ist $4$ m lang und überbrückt waagerecht $3{,}6$ m. Wie groß ist der Höhenunterschied? Runde auf zwei Stellen nach dem Komma. \hfill (P24) \\ \leerfeld[m]
\steil Eine zylinderförmige Dose soll $750\,\text{cm}^3$ fassen und $9$ cm hoch sein. Wie groß ist der Radius? Runde auf zwei Stellen nach dem Komma. \hfill (P22*) \\ \leerfeld[cm]
\steil Ein zylinderförmiges Glas fasst $330$ ml und ist $8$ cm hoch. Wie groß ist der Radius? Runde auf zwei Stellen nach dem Komma. \hfill (P22*) \\ \leerfeld[cm]
\steil $f(x) = x^2 - 8x + 13$ – Nullstellen mit der p-q-Formel? Runde auf zwei Stellen nach dem Komma. \hfill (P25*) \\ x1 = \leerfeld , x2 = \leerfeld
\steil $g(x) = x^2 + 4x - 1$ – Nullstellen mit der p-q-Formel? Runde auf zwei Stellen nach dem Komma. \hfill (P25*) \\ x1 = \leerfeld , x2 = \leerfeld
\end{teile}
\end{aufgabe}

\medskip
\uebersichtskasten{Wurzel: Die Quadratwurzel ist die Umkehrung des Quadrierens – √a ist die Zahl, die mit sich selbst malgenommen a ergibt, und sie ist nie negativ: √144 = 12, weil 12 · 12 = 144. \\ Quadrat und Wurzel heben sich auf – bei negativer Basis erst das Quadrat rechnen, die Wurzel bleibt positiv: √(9²) = 9 und √((−9)²) = √81 = 9. \\ Keine Wurzel aus einer negativen Zahl: √(−81) hat keinen Wert, weil kein Quadrat negativ ist. \\ Näherungswert: Geht die Wurzel nicht auf, liegt sie zwischen zwei Nachbar-Quadratzahlen; der Taschenrechner gibt den Wert, gerundet wird erst am Ende: √30 ≈ 5,48 (zwischen √25 = 5 und √36 = 6).}
